\documentclass[12pt,a4paper]{article}
\usepackage[utf8]{inputenc}
\usepackage[T1]{fontenc}
\usepackage[spanish,es-noshorthands]{babel}
\usepackage{amsmath,amssymb,geometry,fancyhdr,xcolor,listings,tikz,array}
\geometry{margin=2.4cm,headheight=16pt}
\pagestyle{fancy}\fancyhf{}\lhead{UPCT}\rhead{GITS}\lfoot{Prácticas de wxMaxima}\rfoot{\thepage}
\definecolor{azul}{RGB}{23,65,105}
\lstset{basicstyle=\ttfamily\small,backgroundcolor=\color{black!4},frame=single,rulecolor=\color{black!15},breaklines=true,columns=fullflexible}

\begin{document}
\begin{center}{\large Práctica de laboratorio VI}\\[5pt]{\LARGE\bfseries\color{azul} Álgebra lineal: aplicaciones lineales}\end{center}
\noindent Nombre y apellidos: \hrulefill\par\medskip
\noindent\textbf{Objetivos.} Relacionar aplicaciones lineales, matrices y sistemas de ecuaciones mediante ejemplos en $\mathbb R^2$ y $\mathbb R^3$. Ejecutar las celdas con Mayús+Intro y justificar cada resultado. Se emplean vectores por columnas y aritmética exacta.
\section{Linealidad y matriz asociada}
Una aplicación $f:V\to W$ es lineal si $f(\alpha u+\beta v)=\alpha f(u)+\beta f(v)$ para cualesquiera vectores y escalares. En bases ordenadas $\mathcal B=(b_1,\ldots,b_n)$ y $\mathcal D$ de origen y destino, las columnas de $M_{\mathcal D\leftarrow\mathcal B}(f)$ son las coordenadas de $f(b_i)$ en $\mathcal D$. Si ambas bases son canónicas escribimos $M_{\mathcal C}(f)$.

\textbf{Ejemplo.} $f:\mathbb R^2\to\mathbb R^3$, $(x,y)\mapsto(x+y,x-y,2x)$. Calculamos las imágenes de los vectores canónicos y construimos
\[A=\begin{pmatrix}1&1\\1&-1\\2&0\end{pmatrix},\qquad f\begin{pmatrix}x\\y\end{pmatrix}=A\begin{pmatrix}x\\y\end{pmatrix}.\]
Comprueba que la tercera componente es la suma de las dos primeras. Contrasta con $g(x,y)=(x+y,xy)$, que no es lineal.
\begin{lstlisting}
kill(all)$
load("linearalgebra")$
f(x,y):=matrix([x+y],[x-y],[2*x])$
A:addcol(f(1,0),f(0,1));
A.matrix([x],[y]);
f(x,y);
rank(A);
\end{lstlisting}
\section{Núcleo, imagen y fórmula de las dimensiones}
El núcleo es $\ker f=\{u\in V:f(u)=0\}$ y la imagen es $\operatorname{Im} f=\{f(u):u\in V\}$. Ambos son subespacios vectoriales. Para hallar el núcleo resolvemos el sistema homogéneo; para hallar la imagen extraemos una base de las imágenes de los vectores de una base del dominio.
\[\boxed{\dim V=\dim\ker f+\dim\operatorname{Im} f.}\]
\textbf{Ejemplo.} $g:\mathbb R^3\to\mathbb R^2$, $(x,y,z)\mapsto(x+y,y+z)$. El núcleo es el subespacio generado por $\begin{pmatrix}-1\\1\\-1\end{pmatrix}$ y la imagen es $\mathbb R^2$. Se verifica $3=1+2$. Para la aplicación anterior se verifica $2=0+2$; su imagen tiene ecuación $w=u+v$ en las coordenadas $(u,v,w)$ del destino.
\begin{lstlisting}
G:matrix([1,1,0],[0,1,1])$
linsolve([x+y=0,y+z=0],[x,y,z]);
nullspace(G); columnspace(G);
rank(G); nullity(G);
rank(G)+nullity(G);
rank(A)+nullity(A);
\end{lstlisting}
\section{Imagen y antiimagen de un subespacio}
Si $U$ es un subespacio de $V$, su imagen $f(U)$ es el subespacio generado por las imágenes de los vectores de una base de $U$. Si $W_0$ es un subespacio de $W$, su antiimagen es $f^{-1}(W_0)=\{u:f(u)\in W_0\}$. Esta notación no presupone que exista una aplicación inversa.

\textbf{Ejemplo.} Para $g$ del apartado anterior, tomamos $U:\ x-z=0$. Una base viene dada por $\begin{pmatrix}1\\0\\1\end{pmatrix}$ y $\begin{pmatrix}0\\1\\0\end{pmatrix}$. Sus imágenes coinciden: $\begin{pmatrix}1\\1\end{pmatrix}$. Por tanto $g(U)$ tiene ecuación $u-v=0$. Para hallar la antiimagen de la recta $W_0:\ u+v=0$, sustituimos las componentes de $g$ y obtenemos $x+2y+z=0$.
\[\dim f(U)=\dim U-\dim(U\cap\ker f).\]
Además, $\dim f^{-1}(W_0)=\dim\ker f+\dim(W_0\cap\operatorname{Im}f)$, aplicando la fórmula de las dimensiones a la restricción correspondiente.
\begin{lstlisting}
U:matrix([1,0],[0,1],[1,0])$
GU:G.U;
columnspace(GU); rank(GU);
expr:matrix([1,1]).G;
linsolve([x+2*y+z=0],[x,y,z]);
nullspace(expr);
\end{lstlisting}
\section{Inyectividad, sobreyectividad e isomorfismos}
Sea $f:\mathbb R^n\to\mathbb R^m$ con matriz $A$. Es inyectiva si $\ker f=\{0\}$, equivalentemente $\operatorname{rg}A=n$. Es sobreyectiva si $\operatorname{Im}f=\mathbb R^m$, equivalentemente $\operatorname{rg}A=m$. Un isomorfismo es una aplicación lineal biyectiva. Si $m=n$, las tres propiedades son equivalentes.

La aplicación $f$ del primer apartado es inyectiva y no sobreyectiva; $g$ es sobreyectiva y no inyectiva. La aplicación $h:\mathbb R^2\to\mathbb R^2$, $(x,y)\mapsto(2x+y,x+y)$ es un isomorfismo. La aplicación $k:\mathbb R^3\to\mathbb R^3$, $(x,y,z)\mapsto(x+y,y+z,x+2y+z)$ no es inyectiva ni sobreyectiva.
\begin{lstlisting}
H:matrix([2,1],[1,1])$
K:matrix([1,1,0],[0,1,1],[1,2,1])$
for J in [A,G,H,K] do print(matrix_size(J),rank(J),nullity(J))$
\end{lstlisting}
\section{Aplicación inversa}
Si $f$ es un isomorfismo, existe una única aplicación lineal inversa $f^{-1}$ tal que ambas composiciones son las identidades. Su matriz en las bases intercambiadas es $A^{-1}$.

Para $h$ anterior, resolvemos $2x+y=u$, $x+y=v$. Se obtiene $h^{-1}(u,v)=(u-v,2v-u)$. Comprueba las dos composiciones, tanto en forma funcional como matricial.
\begin{lstlisting}
linsolve([2*x+y=u,x+y=v],[x,y]);
Hi:invert(H);
Hi.H; H.Hi;
Hi.matrix([u],[v]);
\end{lstlisting}
\section{Cambio de bases en una aplicación lineal}
Sean $P$ y $Q$ las matrices cuyas columnas son los vectores de las nuevas bases de origen y destino, expresados en las bases antiguas. Si $A$ es la matriz antigua, entonces
\[\boxed{M_{\mathcal D\leftarrow\mathcal B}(f)=Q^{-1}AP.}\]
En efecto, $[u]_{\rm antigua}=P[u]_{\mathcal B}$ y $[f(u)]_{\rm antigua}=Q[f(u)]_{\mathcal D}$. Sustituyendo se deduce la fórmula. Obsérvese que $P$ y $Q$ pueden tener órdenes distintos.

\textbf{Ejemplo.} Para $f:\mathbb R^2\to\mathbb R^3$ del primer apartado elegimos
\[P=\begin{pmatrix}1&1\\0&1\end{pmatrix},\qquad Q=\begin{pmatrix}1&0&1\\0&1&0\\0&0&1\end{pmatrix}.\]
Calculamos la matriz nueva y verificamos la igualdad $QM=AP$. Las columnas de $M$ deben ser las coordenadas de las imágenes de la nueva base de origen en la nueva base de destino.
\begin{lstlisting}
P:matrix([1,1],[0,1])$
Q:matrix([1,0,1],[0,1,0],[0,0,1])$
M:invert(Q).A.P;
Q.M-A.P;
c:matrix([2],[-1])$
Q.M.c; A.P.c;
\end{lstlisting}
\section{Matrices equivalentes y semejantes}
Dos matrices $A,B$ del mismo tamaño son equivalentes si existen matrices invertibles $P,Q$ tales que $B=Q^{-1}AP$. Representan una misma aplicación lineal con elecciones independientes de bases (o aplicaciones identificadas mediante isomorfismos). Son equivalentes si y solo si tienen igual rango: por operaciones elementales de filas y columnas ambas se reducen a la forma normal con un bloque identidad de orden $r$ y ceros en el resto.

Dos matrices cuadradas son semejantes si existe $P$ invertible tal que $B=P^{-1}AP$. Representan un mismo endomorfismo con un cambio simultáneo de base en origen y destino. La semejanza implica equivalencia, pero no recíprocamente. Por ejemplo $I_2$ y $2I_2$ son equivalentes, pero no semejantes: toda matriz semejante a $I_2$ es $I_2$.

\textbf{Ejemplo.} Cambiamos la base del endomorfismo $h$ usando $P$ del apartado anterior. Verificamos la semejanza y la invariancia del rango.
\begin{lstlisting}
Hs:invert(P).H.P;
P.Hs-H.P;
rank(H); rank(Hs);
I:ident(2)$
rank(I); rank(2*I);
invert(P).I.P;
\end{lstlisting}
\section{Interpretación del teorema de Rouché--Frobenius}
El sistema $AX=b$ equivale a buscar la antiimagen de un vector $b$ por la aplicación lineal asociada a $A$. Es compatible si y solo si $b\in\operatorname{Im}f$, es decir,
\[\boxed{\operatorname{rg}A=\operatorname{rg}[A\mid b].}\]
Si $u_0$ es una solución particular, todas las soluciones son $u_0+\ker f$. El conjunto de soluciones es un subespacio afín, no necesariamente vectorial. Si hay $n$ incógnitas y el rango común es $r$, su dimensión afín es $n-r$. Si $r=n$ la solución es única; si $r<n$ hay infinitas soluciones.

\textbf{Ejemplo.} Con $A$ del primer apartado, $b=(3,1,4)$ pertenece a la imagen y tiene una única antiimagen; $b=(3,1,5)$ no pertenece a ella. Con $G$ del segundo apartado, $b=(3,1)$ tiene infinitas antiimágenes. No confundas la antiimagen de un vector con la de un subespacio.
\begin{lstlisting}
b1:matrix([3],[1],[4])$
b2:matrix([3],[1],[5])$
rank(A); rank(addcol(A,b1)); rank(addcol(A,b2));
linsolve([x+y=3,x-y=1,2*x=4],[x,y]);
b3:matrix([3],[1])$
rank(G); rank(addcol(G,b3));
linsolve([x+y=3,y+z=1],[x,y,z]);
G.matrix([3],[0],[1]);
\end{lstlisting}
\section{Ejercicios de trabajo individual}\begin{enumerate}
\item Sea $f:\mathbb R^2\to\mathbb R^3$, $(x,y)\mapsto(2x-y,x+y,3x)$. Obtén la matriz, núcleo, imagen (base y ecuación implícita) y clasifica la aplicación.
\item Sea $g:\mathbb R^3\to\mathbb R^2$, $(x,y,z)\mapsto(x-z,y+z)$. Calcula $g(U)$ para $U:\ x+y=0$ y $g^{-1}(W)$ para $W:\ u-v=0$.
\item Determina para qué valores de $a$ es isomorfismo $h_a:\mathbb R^2\to\mathbb R^2$, $(x,y)\mapsto(x+ay,ax+y)$. Halla la inversa cuando exista.
\item Para $k:\mathbb R^3\to\mathbb R^3$, $(x,y,z)\mapsto(x+y,y+z,z)$, utiliza la base cuyas columnas son las de $Q$ del ejemplo. Calcula la matriz asociada y la aplicación inversa.
\item Para $f$ del ejercicio 1, usa las bases dadas por $P,Q$ del ejemplo. Calcula la matriz asociada y verifica que es equivalente a la canónica.
\item Estudia la compatibilidad del sistema $x+y=2$, $y+z=1$, $x+2y+z=b$ según $b$. Interpreta geométricamente el conjunto de soluciones.
\end{enumerate}
\section{Cuestiones de comprensión}
\begin{enumerate}
\item ¿Puede ser sobreyectiva una aplicación lineal de $\mathbb R^2$ en $\mathbb R^3$? ¿E inyectiva una de $\mathbb R^3$ en $\mathbb R^2$?
\item ¿Por qué la antiimagen de un subespacio es un subespacio aunque la aplicación no sea invertible?
\item ¿Son necesariamente semejantes dos matrices cuadradas equivalentes?
\item ¿Cuándo es vectorial el conjunto no vacío de soluciones de $AX=b$?
\end{enumerate}
\end{document}